\documentclass[11pt,a4paper]{article}
\usepackage[utf8]{inputenc}
\usepackage[T1]{fontenc}
\usepackage[margin=1in]{geometry}
\usepackage{amsmath,amssymb,amsfonts,amsthm}
\usepackage{hyperref}
\usepackage{booktabs}
\usepackage{enumitem}
\usepackage{xcolor}
\usepackage{listings}
\usepackage{tabularx}
\usepackage{../shared/luxcover}

\hypersetup{colorlinks=true,linkcolor=black,citecolor=blue,urlcolor=blue}

\lstset{
  basicstyle=\ttfamily\small,
  breaklines=true,
  frame=single,
  numbers=left,
  numberstyle=\tiny,
  columns=fullflexible
}

\title{Lux Private MPC Deployment\\
\large Fund-only mpcd clusters with no primary network registration}

\author{Lux Research\thanks{Corresponding author: research@lux.network}\\
\textit{Lux Industries, Inc.}\\
\texttt{research@lux.network}}

\date{April 2026}

\begin{document}
\luxcoverpage
\maketitle

\begin{abstract}
This paper specifies the deployment contract for fund-private installations of
\texttt{lux/mpc}, the threshold-signing daemon that backs Lux MPC custody.
The same source ships with a single mode flag, \texttt{--network}, that selects
between two terminations: a \emph{primary} cluster that announces itself to
the Lux M-Chain registrar and anchors audit batches there, and a
\emph{private} cluster that does neither and therefore can run inside an
air-gapped fund VPC. The paper documents the contract, the audit chain that
survives the absence of M-Chain, and the operational procedures a fund's
custody team executes from zero to first signed transaction.
\end{abstract}

\tableofcontents

\section{Why two modes}

A growing number of institutional custodians cannot run their threshold
signing infrastructure on a shared validator network: regulators in their
jurisdictions require every audit byte to be retained by the regulated
entity itself, and disaster-recovery posture forbids a dependency on an
external network for the cryptographic decisions that release client funds.

Lux MPC was built around the Lux M-Chain registrar from the beginning ---
nodes joined a validator-style set, signing decisions were anchored in
batches, and cross-fund operations were possible because every participant
shared the same audit ledger. That model fits well for retail-tier
custody and for the smaller funds that white-label Lux MPC as a service.
It does not fit a \$1B+ AUM fund that runs its custody stack inside a
hardware-protected on-prem facility.

This document specifies the second termination of the same daemon: a
fund deploys mpcd in private mode, every node runs in the fund's own
Kubernetes cluster, the audit log lands on a hardware-WORM volume the
fund controls, and \emph{no} traffic crosses the boundary between the
fund's cluster and any Lux infrastructure unless the fund explicitly
configures it. Same code, same protocols, same operator surface --- the
only thing that changes is the audit destination and the registration
target.

\section{Mode flag and audit dispatcher}

Two flags drive the entire deployment-mode contract:

\begin{lstlisting}
mpcd start \
  --network=primary | private \
  --audit-store=mchain | local-worm | s3-glacier-vault | \
                  azure-immutable | gcs-object-lock | composite
\end{lstlisting}

\subsection{\texttt{--network}}

\texttt{primary} (the default) instructs the bootstrap routine to POST
this node's public key to the registrar URL passed in
\texttt{--registrar-url}; the registrar is responsible for whatever
validator-set bookkeeping the network performs, and the fund need not
care about the protocol details below that HTTP call.

\texttt{private} suppresses the announcement entirely. The node generates
or loads its identity, persists it under \texttt{--keys}, and forms a
cluster strictly with the peers passed via \texttt{--peer}. The flag is
verified at process start: misconfiguration (an unknown mode, or
\texttt{--registrar-url} configured but unreachable in primary mode)
fails closed.

\subsection{\texttt{--audit-store}}

The audit dispatcher decouples \emph{what} we audit from \emph{where} the
log lives. Six dispatcher implementations exist:

\begin{description}
  \item[\texttt{local-worm}] Append-only NDJSON file at
  \texttt{--audit-worm-path}. Replayed on startup; corruption refuses to
  open. Operators mount the path on hardware-WORM media.
  \item[\texttt{mchain}] Anchors \texttt{batchSize} events at a time to a
  Lux M-Chain anchor RPC. Used in primary mode and in private mode when
  the fund has opted into shared anchoring.
  \item[\texttt{s3-glacier-vault}, \texttt{azure-immutable},
  \texttt{gcs-object-lock}] Cloud-WORM backends. Available only when
  the binary is built with the \texttt{cloud-audit} build tag; selecting
  one without the tag fails closed at startup.
  \item[\texttt{composite}] Fans Append out to multiple legs. The
  recommended production posture in primary mode is
  \texttt{composite} with legs \texttt{local-worm,mchain} so a transient
  M-Chain anchor outage cannot lose audit data.
\end{description}

\subsection{Audit chain}

Every event carries a sequence number, the SHA-256 of the previous
event, and the SHA-256 of the canonical bytes of $(\text{prevHash} \,||\,
\text{seq} \,||\, \text{node} \,||\, \text{kind} \,||\, \text{org} \,||\,
\text{wallet} \,||\, \text{timestamp} \,||\, \text{payload})$. Tampering
with any historical event invalidates every subsequent hash; verification
reduces to recomputing the chain. We deliberately do not use
\texttt{json.Marshal} for canonicalisation --- map-iteration ordering in
Go is non-deterministic, which would defeat reproducibility.

\section{Identity bootstrap}

\texttt{pkg/identity/bootstrap.go} owns identity creation and (in
primary mode) registrar announcement. The contract:

\begin{enumerate}
  \item If the keys file exists: load the on-disk Ed25519 keypair, refresh
        the recorded mode, return.
  \item If the keys file does not exist: generate a fresh keypair, write
        atomically with file mode 0600, return.
  \item In primary mode: POST $\{\text{nodeID}, \text{publicKey},
        \text{ts}\}$ to the registrar URL. Failure surfaces; the caller
        decides whether to fall back or fail closed.
\end{enumerate}

In private mode, step 3 is skipped unconditionally. The mode flag is
also recorded in the on-disk JSON so a verifier can detect silent mode
switches between restarts.

\section{Network policy}

The private-mode kustomize overlay
(\texttt{deployments/k8s-private/}) ships a NetworkPolicy that locks
egress to:

\begin{itemize}
  \item Other mpcd pods in the same StatefulSet (P2P transport, port
        9999, and internal API, port 9800).
  \item The fund's postgres and valkey pods (control-plane state).
  \item A pod labelled \texttt{app=fund-hsm-bridge} on port 11111
        (the PKCS\#11 sidecar; operators relabel for their HSM
        provider).
  \item DNS (UDP/TCP 53) only.
\end{itemize}

There is no rule for the public internet, no rule for the Lux M-Chain
RPC, and no rule for an external object store. Funds that opt into a
cloud-WORM target add a rule themselves; the default is silence.

\section{Key ceremony}

The first key the fund's cluster ever produces requires every replica
to participate. Procedure:

\begin{enumerate}[label=\arabic*.]
  \item Verify all replicas report \texttt{healthy} on
        \texttt{/healthz}.
  \item Confirm \texttt{ArePeersReady()} returns true on every node.
  \item Issue a \texttt{POST /keygen} to one node's internal API,
        passing $\{\text{org\_id}, \text{wallet\_id}\}$. The body is
        signed with the operator's intent key.
  \item Wait for the keygen result event (default timeout 60s). The
        node returns the ECDSA and EdDSA public keys plus derived
        addresses for ETH, BTC, SOL.
  \item Verify the result on every other node by querying
        \texttt{/keys/$\langle$wallet\_id$\rangle$}; all replicas must
        return identical metadata.
\end{enumerate}

The ceremony emits a single \texttt{keygen} audit event per node;
operators verify the chain head matches across all replicas before
declaring the wallet provisioned.

\section{Signing ceremony}

Every signing call exercises a $t$-of-$n$ subset. With the default
3-of-3 cluster and threshold $t=2$, two ready nodes form a quorum.
Procedure:

\begin{enumerate}[label=\arabic*.]
  \item Operator submits the transaction payload to the dashboard API.
  \item Approval workflow gates the request on the fund's policy
        (spending limits, time locks, allowlists).
  \item Once approved, the dashboard issues a \texttt{POST /sign} to
        the cluster.
  \item Two or more nodes participate in CGGMP21 (ECDSA) or FROST
        (Ed25519). The signature is returned on the result topic.
  \item A \texttt{sign} audit event is appended to every replica's
        chain. Operators can reconcile the chain across replicas to
        confirm exactly one signing event occurred.
\end{enumerate}

\section{Reshare ceremony}

Reshare rotates key shares without changing the underlying public key.
Used to add a node, remove a node, or change the threshold. Procedure
mirrors the keygen ceremony but uses the LSS protocol; refer to the
mpcd reshare handler for the wire details.

\section{Compliance attestation}

The audit chain produced by mpcd in private mode satisfies the evidence
requirements of:

\begin{itemize}
  \item \textbf{SOC 2 Type II} --- CC6.1, CC6.6, CC7.2 (logical access,
        encryption, change management).
  \item \textbf{ISO 27001:2022} --- A.5.34, A.8.15, A.8.21 (privacy,
        logging, network controls).
  \item \textbf{FIPS 140-2 Level 3} --- when paired with a CloudHSM /
        Azure Key Vault Managed HSM / on-prem PKCS\#11 HSM at L3.
  \item \textbf{CCSS Level III} --- with the additional ceremony witness
        signatures the dashboard captures during keygen and reshare.
\end{itemize}

\section{Honest residual}

The cloud-audit dispatchers (S3 Glacier Vault, Azure Immutable, GCS
Object Lock) compile but return ``not implemented'' at startup. They
are wired into the flag surface so an operator can plan for them; the
SDK paths are gated behind the \texttt{cloud-audit} build tag and the
default release binary does not include them. Funds that want a cloud
WORM target today should use the \texttt{composite} dispatcher with the
\texttt{local-worm} leg pointing at a hardware-WORM-backed PVC and the
\texttt{mchain} leg pointed at their own anchor RPC.

The Terraform modules under \texttt{deployments/terraform/\{aws,gcp,
azure\}/} \texttt{validate} cleanly; only the audit-bucket / VPC
resources are wired end-to-end. EKS / GKE / AKS modules and HSM modules
are stubbed and not yet plan-tested in any cloud account.

\end{document}
